\documentclass[onedown-examples]{subfiles}

\begin{document}
\noindent
\begin{tcolorbox}[colback=red!5!white,colframe=red!75!black,
                  title={Use card symbols like {\setdefaults{colors=b+w}♣,♢,♡,♠} directly as input}]
%\begin{HBox}{Use card symbols like {\setdefaults{colors=b+w}♣,♢,♡,♠} directly as input}
If you want to use the card symbols {\setdefaults{colors=b+w}♣,♢,♡,♠} within your \LaTeX-document,
then put these lines in your preamble. That's all.
\\[0.3em]
\verb|\DeclareUnicodeCharacter{2660}{\Sp}|
\verb|\DeclareUnicodeCharacter{2661}{\He}| \\
\verb|\DeclareUnicodeCharacter{2662}{\Di}|
\verb|\DeclareUnicodeCharacter{2663}{\Cl}| \\
\verb|\DeclareUnicodeCharacter{2664}{\Sp}|
\verb|\DeclareUnicodeCharacter{2665}{\He}| \\
\verb|\DeclareUnicodeCharacter{2666}{\Di}|
\verb|\DeclareUnicodeCharacter{2667}{\Cl}| \\[0.3em]

%\vspace{0.3em}
%\\[0.3em]
If you put them in a file \textbf{\emph{onedown.symbols}}, somewhere where
\LaTeX{} can find it, then the \OneDown package will take care of loading it. An
apropriate place would be in your .texmf directory.
\end{tcolorbox}
%\end{HBox}
\begin{tcolorbox}[colback=red!5!white,colframe=red!75!black,
                  title={To enter these card symbols on your keyboard, use this with Linux}]
Add following lines to your \emph{.Xmodmap} and run \textbf{xmodmap .Xmodmap}
\\[0.3em]
{\ttfamily
keycode 38 = 0x0061 0x0041 0x0061 0x0041 0x1002663 0x1002667 0x00e6 \\
keycode 39 = 0x0073 0x0053 0x0073 0x0053 0x1002662 0x1002666 0x100017f \\
keycode 40 = 0x0064 0x0044 0x0064 0x0044 0x1002661 0x1002665 0x00f0 \\
keycode 41 = 0x0066 0x0046 0x0066 0x0046 0x1002660 0x1002664 0x01f0 \\
keycode 65 = space nobreakspace space nobreakspace nobreakspace NoSymbol \\
}
\tcblower
Now you can enter {\setdefaults{colors=b+w}♣,♢,♡,♠} with \emph{Shift+AltGr} a,s,d,f
\end{tcolorbox}


\end{document}
\endinput
